% Options for packages loaded elsewhere
\PassOptionsToPackage{unicode}{hyperref}
\PassOptionsToPackage{hyphens}{url}
\PassOptionsToPackage{table}{xcolor}
%
\documentclass[
  letterpaper,
]{article}
\usepackage{amsmath,amssymb}
\usepackage{iftex}
\usepackage[T1]{fontenc}
\usepackage[utf8]{inputenc}
\usepackage{textcomp}
\usepackage{lmodern}
\ifPDFTeX\else
  % xetex/luatex font selection
\fi
% Use upquote if available, for straight quotes in verbatim environments
\IfFileExists{upquote.sty}{\usepackage{upquote}}{}
\IfFileExists{microtype.sty}{% use microtype if available
  \usepackage[]{microtype}
  \UseMicrotypeSet[protrusion]{basicmath} % disable protrusion for tt fonts
}{}
\makeatletter
\@ifundefined{KOMAClassName}{% if non-KOMA class
  \IfFileExists{parskip.sty}{%
    \usepackage{parskip}
  }{% else
    \setlength{\parindent}{0pt}
    \setlength{\parskip}{6pt plus 2pt minus 1pt}}
}{% if KOMA class
  \KOMAoptions{parskip=half}}
\makeatother
\usepackage{xcolor}
\usepackage{color}
\usepackage{fancyvrb}
\newcommand{\VerbBar}{|}
\newcommand{\VERB}{\Verb[commandchars=\\\{\}]}
\DefineVerbatimEnvironment{Highlighting}{Verbatim}{commandchars=\\\{\}}
% Add ',fontsize=\small' for more characters per line
\newenvironment{Shaded}{}{}
\newcommand{\AlertTok}[1]{\textcolor[rgb]{1.00,0.00,0.00}{\textbf{#1}}}
\newcommand{\AnnotationTok}[1]{\textcolor[rgb]{0.38,0.63,0.69}{\textbf{\textit{#1}}}}
\newcommand{\AttributeTok}[1]{\textcolor[rgb]{0.49,0.56,0.16}{#1}}
\newcommand{\BaseNTok}[1]{\textcolor[rgb]{0.25,0.63,0.44}{#1}}
\newcommand{\BuiltInTok}[1]{\textcolor[rgb]{0.00,0.50,0.00}{#1}}
\newcommand{\CharTok}[1]{\textcolor[rgb]{0.25,0.44,0.63}{#1}}
\newcommand{\CommentTok}[1]{\textcolor[rgb]{0.38,0.63,0.69}{\textit{#1}}}
\newcommand{\CommentVarTok}[1]{\textcolor[rgb]{0.38,0.63,0.69}{\textbf{\textit{#1}}}}
\newcommand{\ConstantTok}[1]{\textcolor[rgb]{0.53,0.00,0.00}{#1}}
\newcommand{\ControlFlowTok}[1]{\textcolor[rgb]{0.00,0.44,0.13}{\textbf{#1}}}
\newcommand{\DataTypeTok}[1]{\textcolor[rgb]{0.56,0.13,0.00}{#1}}
\newcommand{\DecValTok}[1]{\textcolor[rgb]{0.25,0.63,0.44}{#1}}
\newcommand{\DocumentationTok}[1]{\textcolor[rgb]{0.73,0.13,0.13}{\textit{#1}}}
\newcommand{\ErrorTok}[1]{\textcolor[rgb]{1.00,0.00,0.00}{\textbf{#1}}}
\newcommand{\ExtensionTok}[1]{#1}
\newcommand{\FloatTok}[1]{\textcolor[rgb]{0.25,0.63,0.44}{#1}}
\newcommand{\FunctionTok}[1]{\textcolor[rgb]{0.02,0.16,0.49}{#1}}
\newcommand{\ImportTok}[1]{\textcolor[rgb]{0.00,0.50,0.00}{\textbf{#1}}}
\newcommand{\InformationTok}[1]{\textcolor[rgb]{0.38,0.63,0.69}{\textbf{\textit{#1}}}}
\newcommand{\KeywordTok}[1]{\textcolor[rgb]{0.00,0.44,0.13}{\textbf{#1}}}
\newcommand{\NormalTok}[1]{#1}
\newcommand{\OperatorTok}[1]{\textcolor[rgb]{0.40,0.40,0.40}{#1}}
\newcommand{\OtherTok}[1]{\textcolor[rgb]{0.00,0.44,0.13}{#1}}
\newcommand{\PreprocessorTok}[1]{\textcolor[rgb]{0.74,0.48,0.00}{#1}}
\newcommand{\RegionMarkerTok}[1]{#1}
\newcommand{\SpecialCharTok}[1]{\textcolor[rgb]{0.25,0.44,0.63}{#1}}
\newcommand{\SpecialStringTok}[1]{\textcolor[rgb]{0.73,0.40,0.53}{#1}}
\newcommand{\StringTok}[1]{\textcolor[rgb]{0.25,0.44,0.63}{#1}}
\newcommand{\VariableTok}[1]{\textcolor[rgb]{0.10,0.09,0.49}{#1}}
\newcommand{\VerbatimStringTok}[1]{\textcolor[rgb]{0.25,0.44,0.63}{#1}}
\newcommand{\WarningTok}[1]{\textcolor[rgb]{0.38,0.63,0.69}{\textbf{\textit{#1}}}}
\usepackage{longtable,booktabs,array}
\usepackage{calc} % for calculating minipage widths
% Correct order of tables after \paragraph or \subparagraph
\usepackage{etoolbox}
\makeatletter
\patchcmd\longtable{\par}{\if@noskipsec\mbox{}\fi\par}{}{}
\makeatother
% Allow footnotes in longtable head/foot
\IfFileExists{footnotehyper.sty}{\usepackage{footnotehyper}}{\usepackage{footnote}}
\makesavenoteenv{longtable}
\setlength{\emergencystretch}{3em} % prevent overfull lines
\providecommand{\tightlist}{%
  \setlength{\itemsep}{0pt}\setlength{\parskip}{0pt}}
\setcounter{secnumdepth}{5}
% Self-contained presentation rules embedded in every generated LaTeX file.
\usepackage{iftex}
\ifPDFTeX\else\PackageError{SSE-ASAS}{Compile with pdfLaTeX only}{}\fi
\usepackage[T1]{fontenc}
\usepackage[utf8]{inputenc}
\usepackage{lmodern}
\usepackage{microtype}
\usepackage{geometry}
\geometry{margin=0.93in,headheight=16pt}
\usepackage{xcolor} % [table] passed at top of file (option-clash fix)
\definecolor{ink}{HTML}{15283B}
\definecolor{brand}{HTML}{075B75}
\definecolor{brandlight}{HTML}{EAF5F8}
\definecolor{warm}{HTML}{A15C1D}
\definecolor{muted}{HTML}{536577}
\usepackage{fancyhdr}
\usepackage{titlesec}
\usepackage{enumitem}
\usepackage{booktabs}
\usepackage{longtable}
\usepackage{array}
\usepackage{tabularx}
\usepackage{tcolorbox}
\tcbuselibrary{breakable,skins}
\usepackage{listings}
\usepackage{fancyvrb}
\usepackage{hyperref}
\hypersetup{colorlinks=true,linkcolor=brand,urlcolor=brand,citecolor=brand}
\urlstyle{same}
\usepackage{lastpage}
\setlength{\emergencystretch}{3em}
\setlength{\parskip}{0.45em}
\setlength{\parindent}{0pt}
\setlist[itemize]{leftmargin=*,topsep=3pt,itemsep=2pt}
\setlist[enumerate]{leftmargin=*,topsep=3pt,itemsep=2pt}
\titleformat{\section}{\Large\bfseries\color{ink}}{\thesection}{0.75em}{}[\vspace{-2pt}\color{brand}\titlerule]
\titleformat{\subsection}{\large\bfseries\color{brand}}{\thesubsection}{0.7em}{}
\titleformat{\subsubsection}{\normalsize\bfseries\color{ink}}{\thesubsubsection}{0.7em}{}
\pagestyle{fancy}
\fancyhf{}
\fancyhead[L]{\small\sffamily\textcolor{brand}{SOFTWARE SECURITY ENGINEERING}}
\fancyhead[R]{\small\sffamily\textcolor{muted}{TECHNICAL SERIES}}
\fancyfoot[L]{\scriptsize\sffamily Version 1.0 \textbullet\ 14 October 2026}
\fancyfoot[R]{\scriptsize\sffamily Page \thepage\ of \pageref{LastPage}}
\renewcommand{\headrulewidth}{0.5pt}
\renewcommand{\footrulewidth}{0.2pt}
\setcounter{tocdepth}{2}
\setcounter{secnumdepth}{3}
\lstdefinestyle{sse}{basicstyle=\ttfamily\footnotesize,frame=single,framerule=0.4pt,rulecolor=\color{brand!40},backgroundcolor=\color{brandlight!45},numbers=left,numberstyle=\scriptsize\color{muted},numbersep=6pt,breaklines=true,breakatwhitespace=false,tabsize=2,showstringspaces=false,columns=fullflexible,keepspaces=true,captionpos=b,aboveskip=9pt,belowskip=9pt}
\lstset{style=sse}
\newtcolorbox{sourceconcept}[1][]{enhanced,breakable,colback=brandlight,colframe=brand!65!black,title={ORIGINAL CONCEPT},fonttitle=\bfseries\sffamily,#1}
\newtcolorbox{modernextension}[1][]{enhanced,breakable,colback=white,colframe=warm!85!black,title={MODERN EXTENSION},fonttitle=\bfseries\sffamily,#1}
\newtcolorbox{fieldwork}[1][]{enhanced,breakable,colback=white,colframe=ink!55,title={PRACTICAL APPLICATION},fonttitle=\bfseries\sffamily,#1}
\AtBeginDocument{\color{ink}}
\makeatletter
\renewcommand{\maketitle}{%
 \begin{titlepage}\thispagestyle{empty}
 \vspace*{1.5cm}{\color{brand}\rule{\textwidth}{3pt}}\par\vspace{1cm}
 {\sffamily\bfseries\color{muted}\large SOFTWARE SECURITY ENGINEERING\\ APPLICATION SECURITY ASSESSMENT SERIES}\par\vspace{1.05cm}
 {\sffamily\bfseries\color{ink}\fontsize{27}{34}\selectfont\@title\par}\vspace{0.8cm}
 {\sffamily\large\color{brand}\@author\par}\vfill
 \begin{tcolorbox}[colback=brandlight,colframe=brand!55,boxrule=0.5pt]
 \small\textbf{Editorial policy}\quad Original teaching synthesis with book-grounded concepts distinguished from contemporary practices. Modern platform behavior is version-sensitive; validate against linked vendor documentation before production use.
 \end{tcolorbox}
 {\color{brand}\rule{\textwidth}{2pt}}\par\vspace{0.3cm}
 {\sffamily\small Version 1.0 \quad|\quad October 2026 \quad|\quad pdfLaTeX}\end{titlepage}\setcounter{page}{1}}
\makeatother
\ifLuaTeX
  \usepackage{selnolig}  % disable illegal ligatures
\fi
\usepackage{bookmark}
\IfFileExists{xurl.sty}{\usepackage{xurl}}{} % add URL line breaks if available
\urlstyle{same}
\hypersetup{
  pdftitle={Manual Source Code Auditing},
  pdfauthor={Independent Technical Study Edition},
  hidelinks,
  pdfcreator={LaTeX via pandoc}}

\title{Manual Source Code Auditing}
\author{Independent Technical Study Edition}
\date{}

\begin{document}
\maketitle

{
\setcounter{tocdepth}{3}
\tableofcontents
}
\section{Mission and boundaries}\label{mission-and-boundaries}

Manual auditing is a disciplined investigation of whether the actual
code and the architecture satisfy specified security properties. It
complements SAST: tools find recognized local patterns and some
interprocedural paths; reviewers uncover omitted checks, business-logic
assumptions, incomplete threat models, and subtle state transitions.

\textbf{Original foundation:} Dowd, McDonald, and Schuh, Chapters 1--4,
especially design review, operational review, and the application review
process. Their iterative \emph{Plan--Work--Reflect} approach,
code-comprehension strategies, candidate-point strategies, and
design-generalization strategies provide this guide's organizing method.
\textbf{Modern adaptation:} pull-request reviews, repository inventory,
CI evidence, and CodeQL findings become new sources of audit leads.

\section{Preparation and scope
control}\label{preparation-and-scope-control}

\subsection{Define the review
contract}\label{define-the-review-contract}

Record: system purpose; in-scope repos and commit SHAs; branches;
language/framework versions; deployment assumptions; architecture and
data-classification diagrams; credentials and access level; time budget;
forbidden testing activities; contacts; output format; and severity
convention. If runtime testing is not authorized, stay with static
reasoning or an explicitly approved isolated environment.

A useful inventory joins entry point to business capability,
authorization policy, implementation, data store, and downstream
consumer. Unmapped routes and scheduled jobs are first-class findings
because they can bypass security assumptions made by HTTP endpoints.

\begin{longtable}[]{@{}ll@{}}
\toprule\noalign{}
Inventory field & Example \\
\midrule\noalign{}
\endhead
\bottomrule\noalign{}
\endlastfoot
Entry point & \texttt{POST\ /api/invoices/\{id\}/approve} \\
Input & Invoice identifier; requester context \\
Decision & Owner permission and distinct approver \\
Sink & Database status update and event publication \\
Trust boundary & Untrusted request to privileged accounting database \\
Evidence & Route source, authorization helper, integration test \\
\end{longtable}

\subsection{Construct a review
hypothesis}\label{construct-a-review-hypothesis}

Examples: (a) all invoice queries include tenant scope; (b) C parsers
always validate \texttt{length} before allocation and copy; (c)
request-derived file paths cannot escape the allowed directory; (d)
logout invalidates sensitive session state. Form a hypothesis that can
be disproved. ``Find vulnerabilities'' is too vague to guide review.

\section{Iterative auditing
procedure}\label{iterative-auditing-procedure}

\subsection{Plan--Work--Reflect}\label{planworkreflect}

\textbf{Plan:} select a security objective, candidate module, and trace
strategy. Set an explicit stopping condition: every entry point in a
subsystem traced to an authorization decision, or every call to an
unsafe file API reviewed.

\textbf{Work:} build a data-flow and control-flow picture, inspect
relevant code, record a call chain, test assumptions against nearby code
and tests, and gather reproducible evidence.

\textbf{Reflect:} assess coverage, update architecture understanding,
choose a new hypothesis, mark assumptions requiring owner confirmation,
and prevent duplicate issues. The next cycle should use knowledge gained
from previous cycles.

\begin{sourceconcept}
Source-based review strategies are complementary: **code comprehension** reconstructs intended behavior; **candidate-point analysis** starts from suspicious sources/sinks/APIs; **design generalization** compares repeated implementations to discover policy inconsistencies. An audit is not simply a search for dangerous function names.
\end{sourceconcept}

\subsection{Top-down, bottom-up, and hybrid
traversals}\label{top-down-bottom-up-and-hybrid-traversals}

\begin{itemize}
\tightlist
\item
  \textbf{Top-down:} threat model \(\to\) architecture \(\to\) component
  \(\to\) function \(\to\) branch. Best for authorization, trust
  assumptions, and business logic.
\item
  \textbf{Bottom-up:} risky sink \(\to\) callers \(\to\) input
  provenance and authorization. Best for parser, memory, logging,
  cryptographic, and interpreter misuse.
\item
  \textbf{Hybrid:} inventory entry points and invariants first; use
  targeted sink searches, then return to the design model when paths
  contradict assumptions.
\end{itemize}

\section{Focused audit playbooks}\label{focused-audit-playbooks}

\subsection{Authentication and object-level
authorization}\label{authentication-and-object-level-authorization}

Trace all routes to protected objects, including bulk actions,
background jobs, GraphQL resolvers, and exports. Distinguish
authentication from object scope; determine whether tenant and ownership
checks occur on the authoritative data path. Search for alternate
operations that update the same record without equivalent guards.

\begin{Shaded}
\begin{Highlighting}[]
\CommentTok{\# Problem: ID selected first, no ownership check.}
\NormalTok{invoice }\OperatorTok{=}\NormalTok{ db.get\_invoice(request.args[}\StringTok{"id"}\NormalTok{])}
\ControlFlowTok{return}\NormalTok{ serialize(invoice)}

\CommentTok{\# Better: scope at query boundary, then handle absence uniformly.}
\NormalTok{invoice }\OperatorTok{=}\NormalTok{ db.get\_invoice\_for\_tenant(}
\NormalTok{    request.args[}\StringTok{"id"}\NormalTok{], principal.tenant\_id)}
\ControlFlowTok{if}\NormalTok{ invoice }\KeywordTok{is} \VariableTok{None}\NormalTok{:}
    \ControlFlowTok{raise}\NormalTok{ NotFound()}
\ControlFlowTok{return}\NormalTok{ serialize(invoice)}
\end{Highlighting}
\end{Shaded}

\textbf{Trace questions:} who supplied the ID; where does identity
originate; can the object change owners; can a privileged service path
bypass the check; does bulk export use the same authorization policy;
can a cache return a cross-tenant result?

\subsection{Native memory and integer
arithmetic}\label{native-memory-and-integer-arithmetic}

Use a \emph{size ledger}: enumerate every buffer, allocation size,
arithmetic conversion, bound, and copy in a call path. Evaluate integer
promotions, signedness, truncation, allocation failure, overlapping
buffers, and object lifetime. Favor a forward trace from untrusted
lengths and a backward trace from memory-writing operations.

\begin{longtable}[]{@{}
  >{\raggedright\arraybackslash}p{(\columnwidth - 6\tabcolsep) * \real{0.2500}}
  >{\raggedright\arraybackslash}p{(\columnwidth - 6\tabcolsep) * \real{0.2500}}
  >{\raggedright\arraybackslash}p{(\columnwidth - 6\tabcolsep) * \real{0.2500}}
  >{\raggedright\arraybackslash}p{(\columnwidth - 6\tabcolsep) * \real{0.2500}}@{}}
\toprule\noalign{}
\begin{minipage}[b]{\linewidth}\raggedright
Variable
\end{minipage} & \begin{minipage}[b]{\linewidth}\raggedright
Source / type
\end{minipage} & \begin{minipage}[b]{\linewidth}\raggedright
Invariant
\end{minipage} & \begin{minipage}[b]{\linewidth}\raggedright
Inspection point
\end{minipage} \\
\midrule\noalign{}
\endhead
\bottomrule\noalign{}
\endlastfoot
\texttt{count} & Header / unsigned 32-bit &
\texttt{count\ \textless{}=\ hard\_limit} & After parse \\
\texttt{bytes} & \texttt{count\ *\ record\_size} & Product representable
& Before allocation \\
\texttt{buf} & Heap allocation & Allocation succeeded & Before every
write \\
\texttt{offset} & Loop arithmetic &
\texttt{offset\ +\ width\ \textless{}=\ bytes} & Each iteration \\
\end{longtable}

\subsection{Command, SQL, path, and markup
injection}\label{command-sql-path-and-markup-injection}

Document the interpreter. For each externally influenced string,
identify all transformations, the exact sink, the surrounding grammar,
and whether safety relies on blacklisting. Parameterized SQL, a no-shell
process API, context-specific HTML encoding, and canonical path
confinement solve different problems; one cannot stand in for another.

\subsection{Concurrency and TOCTOU}\label{concurrency-and-toctou}

List shared state, locks, transactions, cache lifetimes, file
operations, and the time between check and use. Challenge the assumption
that two requests execute serially. Trace async jobs and webhooks that
independently alter the same record. A deterministic multi-request unit
test with a barrier is often better evidence than a timing-sensitive
manual race demo.

\subsection{Secrets, authentication tokens, and diagnostic
surfaces}\label{secrets-authentication-tokens-and-diagnostic-surfaces}

Examine credentials in source, build scripts, artifacts, error handlers,
telemetry, and local test fixtures. For a real exposed credential, the
first priority is to contain/revoke/rotate it under the incident
process; deleting a line from the latest commit alone does not make the
secret unexposed. Handle artifacts carefully and avoid copying secret
values into findings.

\section{Search, trace, and validation
technique}\label{search-trace-and-validation-technique}

Start from a small query set, then validate each candidate in context.
For native code, search allocation and copy sites; for web code, find
database callers and response construction; for concurrency, identify
mutable shared stores.

\begin{Shaded}
\begin{Highlighting}[]
\CommentTok{\# Examples run only against repositories you are authorized to assess.}
\ExtensionTok{rg} \AttributeTok{{-}n} \StringTok{\textquotesingle{}execute\textbackslash{}(|cursor\textbackslash{}.execute|exec\textbackslash{}(|subprocess|render\_template\textquotesingle{}}\NormalTok{ .}
\ExtensionTok{rg} \AttributeTok{{-}n} \StringTok{\textquotesingle{}memcpy|strcpy|sprintf|malloc|realloc|strncat\textquotesingle{}} \AttributeTok{{-}{-}glob} \StringTok{\textquotesingle{}*.\{c,h,cpp\}\textquotesingle{}}\NormalTok{ .}
\ExtensionTok{rg} \AttributeTok{{-}n} \StringTok{\textquotesingle{}tenant\_id|organization\_id|authorize|permission\textquotesingle{}}\NormalTok{ src test}
\end{Highlighting}
\end{Shaded}

Static grep results are candidates, not validated findings. Inspect
macro expansion, framework-provided escaping, parameter semantics,
symbol aliasing, invocation conditions, and reachable paths. Conversely,
absence of a matching string does not rule out a vulnerability.

\section{Evidence packet and stopping
rules}\label{evidence-packet-and-stopping-rules}

\textbf{Evidence packet:} immutable source reference; boundary/actor;
violated invariant; minimal trace (source \(\to\) transformation \(\to\)
sink); conditions; actual-versus-expected result; severity argument;
remediation options; negative regression test; reviewer confidence.

\textbf{Stop when:} all planned paths have dispositions; critical sinks
have full provenance; competing paths to an asset have been considered;
unverified assumptions are tagged; observed results are reproducible;
scope gaps are declared. Claim ``not observed within scope,'' not ``no
vulnerabilities exist.''

\begin{fieldwork}
**Audit simulation:** inspect a two-endpoint order service that supports `GET /orders/{id}` and `POST /orders/{id}/refund`. Produce an entry-point map, one authorization invariant per endpoint, a counterexample involving cross-tenant access, and two tests that must fail against an intentionally flawed version. Provide precise source locations and scope limitations.
\end{fieldwork}

\section{Quality rubric}\label{quality-rubric}

\begin{longtable}[]{@{}
  >{\raggedright\arraybackslash}p{(\columnwidth - 4\tabcolsep) * \real{0.3333}}
  >{\raggedright\arraybackslash}p{(\columnwidth - 4\tabcolsep) * \real{0.3333}}
  >{\raggedright\arraybackslash}p{(\columnwidth - 4\tabcolsep) * \real{0.3333}}@{}}
\toprule\noalign{}
\begin{minipage}[b]{\linewidth}\raggedright
Dimension
\end{minipage} & \begin{minipage}[b]{\linewidth}\raggedright
Weak report
\end{minipage} & \begin{minipage}[b]{\linewidth}\raggedright
Strong report
\end{minipage} \\
\midrule\noalign{}
\endhead
\bottomrule\noalign{}
\endlastfoot
Root cause & ``Missing security'' & Missing object-scope predicate on a
concrete call path \\
Evidence & Scanner screenshot only & Commit SHA, trace,
input/precondition, observed behavior \\
Remediation & ``Sanitize'' & Enforce tenant predicate at canonical data
access and test it \\
Confidence & Not stated & Observed / inferred / needs runtime
confirmation \\
Completeness & No mention of adjacent paths & Explains nearby bypasses
and review boundaries \\
\end{longtable}

\section{Annex A: repeatable auditor
worksheets}\label{annex-a-repeatable-auditor-worksheets}

\subsection{Authorization trace
worksheet}\label{authorization-trace-worksheet}

Fill the following for each sensitive endpoint: \texttt{route},
\texttt{HTTP\ method\ or\ background\ trigger},
\texttt{actor\ identity\ source}, \texttt{requested\ object},
\texttt{ownership\ lookup}, \texttt{role/action\ decision},
\texttt{transactional\ state\ check}, \texttt{return\ data},
\texttt{error\ behavior}, \texttt{tests}. Use a separate row for list,
search, export, report generation, and delegated service calls. If a
field is unknown, do not replace it with ``secure''. Open a scoped
evidence request with a named owner.

\textbf{High-value edge cases:} unauthenticated request; disabled
principal; expired delegation; deleted resource; unknown tenant;
legitimate tenant with no action permission; staff support
impersonation; valid foreign object; stale cache entry; downgraded
subscription; retry of an already-finalized state transition. Most
require only synthetic test data and a normal application test harness.

\subsection{Backward trace from a dangerous
sink}\label{backward-trace-from-a-dangerous-sink}

Start with \texttt{subprocess} or \texttt{exec} at the lowest layer.
Record the executable path, argument-vector construction, shell usage,
environment variables, current working directory, file paths, timeouts,
identity, and resulting side effects. Follow each argument backward to
its source, flag transformations, and examine whether the operation
really needs to invoke an external process. Check success/failure
cleanup and diagnostic output. A no-shell call avoids shell parsing but
an attacker may still influence target-program options. A secure wrapper
validates arguments against the target program's \emph{own} semantics.

For a database sink, record driver API, prepared-statement usage,
dynamic identifiers, authorization scoping, least-privilege DB role,
transaction boundaries, and error behavior. Parameterized values do not
authorize a row; authorization scoping is an independent requirement.

\subsection{Cross-file C/C++ audit
ledger}\label{cross-file-cc-audit-ledger}

\begin{longtable}[]{@{}
  >{\raggedright\arraybackslash}p{(\columnwidth - 4\tabcolsep) * \real{0.3333}}
  >{\raggedright\arraybackslash}p{(\columnwidth - 4\tabcolsep) * \real{0.3333}}
  >{\raggedright\arraybackslash}p{(\columnwidth - 4\tabcolsep) * \real{0.3333}}@{}}
\toprule\noalign{}
\begin{minipage}[b]{\linewidth}\raggedright
Operation
\end{minipage} & \begin{minipage}[b]{\linewidth}\raggedright
Reviewer question
\end{minipage} & \begin{minipage}[b]{\linewidth}\raggedright
Evidence type
\end{minipage} \\
\midrule\noalign{}
\endhead
\bottomrule\noalign{}
\endlastfoot
Parsing & Could count/length be negative or truncated? & Input type and
range proof \\
Allocation & Can \texttt{n\ *\ size} overflow or return null? & Guard
and tested boundary \\
Conversion & Does unit or signedness change? & Cast and calling
convention \\
Copy & Are both source and destination lengths known? & Size ledger and
contract \\
Ownership & Can an alias outlive or free the storage? & Lifetime and
cleanup trace \\
Concurrency & Can another thread mutate buffer/state? & Synchronization
and schedule tests \\
\end{longtable}

Use the above as a \emph{call-chain audit} rather than treating each
statement independently. Review the native boundary exposed to Python,
Java, web servers, and parsers where data arrives through
foreign-language interfaces. Monitor unusual invariants: pointer
arithmetic across allocations, indexes based on previous frames,
variable-sized structures, and library calls that interpret NUL
termination differently.

\subsection{Documentation template: audit cycle
record}\label{documentation-template-audit-cycle-record}

\begin{Shaded}
\begin{Highlighting}[]
\NormalTok{Cycle ID:          AUD{-}YYYY{-}NNN}
\NormalTok{Repository/SHA:    ...}
\NormalTok{Security property: ...}
\NormalTok{Working hypothesis:...}
\NormalTok{Chosen strategy:   top{-}down / candidate point / hybrid}
\NormalTok{Entry point(s):    ...}
\NormalTok{Sensitive sinks:   ...}
\NormalTok{Relevant call path:...}
\NormalTok{Observed evidence: ...}
\NormalTok{Open assumptions:  ...}
\NormalTok{Contradictions:     ...}
\NormalTok{Coverage achieved: ...}
\NormalTok{Next{-}cycle action: ...}
\NormalTok{Reviewer/date:     ...}
\end{Highlighting}
\end{Shaded}

\textbf{Quality review:} Can another engineer reconstruct the call
chain? Were successful and exceptional execution paths examined? Was a
path deemed impossible because of code evidence or merely because a
developer said it should be impossible? Are there alternate
implementations of the same security feature? Is a compiler/build flag
necessary for the conclusion?

\section{Annex B: triage without collapsing
categories}\label{annex-b-triage-without-collapsing-categories}

A manual audit can invalidate a scanner report without proving a program
safe. Example: a tool flags \texttt{cursor.execute(query)} as SQL
injection. The reviewer discovers that \texttt{query} is a literal
selected from a static enum and the only user field is passed as a bound
parameter. The \textbf{specific scanner path} is a false positive if the
data model is documented accurately. But a separate endpoint using the
same database function may have missing object authorization. Record
both conclusions independently.

A mitigated finding likewise differs from a false positive: the unsafe
pattern may remain, but a real control makes the relevant attack path
unreachable under validated assumptions. Record where the mitigation
sits, who controls it, how it is monitored, and what would happen if it
were removed. When uncertainty remains, assign a follow-up verification
task rather than claiming decisive evidence.

\subsection{Audit reporting example: from hypothesis to
proof}\label{audit-reporting-example-from-hypothesis-to-proof}

\textbf{Hypothesis:} request-supplied \texttt{invoice\_id} can select
another tenant's record. \textbf{Trace:} router -\textgreater{} service
-\textgreater{} repository \texttt{findById} -\textgreater{} response.
\textbf{Observation:} test account A receives a fixture record owned by
B. \textbf{Counterfactual:} expected lookup includes \texttt{tenant\_id}
from validated principal and returns not found for foreign objects.
\textbf{Cause:} missing resource scoping at an authoritative data access
point. \textbf{Patch:} centralize \texttt{findForTenant}; restrict
generic \texttt{findById} to trusted internal use. \textbf{Regressions:}
read, bulk export, and notification generation reject cross-tenant
access. \textbf{Limit:} jobs triggered outside HTTP still require
separate inspection.

The narrative identifies violation, evidence, cause, remediation, and
coverage limit. It is significantly more actionable than ``critical
authorization bug in endpoint''.

\section{Source and modern extension
map}\label{source-and-modern-extension-map}

\textbf{Book-grounded:} Dowd et al., Chapter 2 (design review and threat
modeling), Chapter 3 (operational review), Chapter 4 (application
review, strategies, toolboxes), Chapters 5--13 (native/platform/state
flaws), Chapters 16--18 (network and Web application review). Chess \&
West, Chapter 3 (static analysis and review workflow).

\textbf{Modern extension:} OWASP ASVS
\url{https://owasp.org/projects/asvs}; GitHub code scanning alerts
\url{https://docs.github.com/en/code-security/concepts/code-scanning/code-scanning-alerts};
NIST SSDF \url{https://csrc.nist.gov/pubs/sp/800/218/final}.

\end{document}
